defined by @@M0@@ for @@M1@@, @@M2@@. It factors through the inclusion of @@M3@@ in @@M4@@, which shows that the homomorphism from @@M5@@ to @@M6@@ is injective, and will allow us to identify the former with a subgroup of the latter; similarly, we shall regard @@M7@@ as contained in @@M8@@.

Notice that @@M9@@ is precisely the augmentation of @@M10@@.

\sgaStableTarget{sga6.e05.para.1-5}1.5. We embed @@M11@@ in @@M12@@ by sending @@M13@@ to the @@M14@@-th symmetric function in the @@M15@@; letting the symmetric group @@M16@@ act by permutation of the @@M17@@, and using the theorem on symmetric polynomials and (1.3.ii)), we see that the resulting homomorphism from @@M18@@ to @@M19@@ is injective, which will again allow us to identify the former with a subgroup of the latter:
@@M20@@
We denote by
@@M21@@
the homomorphism defined by @@M22@@, and put
@@M23@@
The symmetric group permutes the morphisms @@M24@@, and therefore leaves the morphism @@M25@@ invariant, which consequently takes its values in
